\scanpage{54}
D'autre part, si $E\in\operatorname{Ob}\unsure{\mathcal M|M}^{\deleted{\unsure{+|\ast}}}(X)$, \deleted{\unsure{p. ex. simple|simple}}, on trouve
% slanted
\margin{$\nearrow$ ``Conjectures de Tate'' $\searrow$}
\[
\left\{\begin{aligned}
\operatorname{cl}(E)\in M^+(X)\quad\Longleftrightarrow\quad
&\forall\,x\ \text{pt fini de }X,\ \text{le polynôme}\\
&\text{caractéristique }\deleted{\text{de Frobenius}}\\
&\operatorname{Frob}_x^{-1}\ \text{dans }T_\ell(E)(x)\text{ est}\\
&\text{à coefficients }\text{\uline{entiers}}\\
&[\,\text{i.e. les valeurs propres de}\\
&\text{Frobenius sont des}\\
&\text{\uline{entiers} algébriques}\,]
\end{aligned}\right.
\]

Par ailleurs, pour \deleted{\unsure{un élément, pour|décrit, pour} $\unsure{\zeta^\alpha}\in M_\alpha(X)$} $E$ ``homogène'' de degré $\alpha$, pour exprimer si $E\in\operatorname{Ob}D_i(M(X))$, où $i\geq 0$ \deleted{on peut, \unsure{par}}
\inserted{\deleted{$\operatorname{cl}(E)\in\unsure{\Sigma^i|\sum_i}$} \deleted{$\unsure{\zeta^\alpha}\in D_\alpha M(X)$}
$E\in\operatorname{Ob}D_i\mathcal M(X)$ \unsure{\}} i.e. $\operatorname{cl}(E)\in D_i M(X)\Longleftrightarrow$ Si} on est $\Downarrow$

$[\,1^\circ)$ $i=0$. Alors
\[
\begin{aligned}
E\in\operatorname{Ob}D_0(M(X))\Longleftrightarrow\quad &E\text{ est de degré \uline{pair} }\alpha=2\beta\\
&\text{et }E\unsure{\zeta|\xi}^{-\beta}=F\quad(\text{qui est de}\\
&\text{degré }0)\text{ est }\in\operatorname{Ob}M^+(X)\\
&\text{i.e. si }\rho_{\unsure{\nu|i}}(x)\text{ sont les}\\
&\text{valeurs propres de }\operatorname{Frob}_x^{-1}\text{ dans}\\
&T_\ell(M(x)),\ \text{alors les }\rho_\nu(x)/N(x)^\beta\\
&\text{sont des entiers algébriques}\,]
\end{aligned}
\]
\margin{N.B. Ce sont bien des racines de l'unité [leur produit étant $\pm1$ !]}

$2^\circ)$ $i\geq 1$. Alors il existe une plus petite puissance $\zeta^\beta$ de $\zeta$ telle que $E\zeta^{-\beta}\deleted{\unsure{\in\operatorname{Ob} M}}\inserted{{}=E(+\beta)}$ soit de degré $i-1$ ou $i$ \unsure{[|(}suivant la parité de $\alpha\,]$. On veut que \deleted{\unsure{cet}} $E(\unsure{+\beta|\beta})$ soit $\in\operatorname{Ob} M^+(X)$, i.e. les $\rho_\nu(x)/N(x)^\beta$ doivent être des entiers alg.
% archive sheet number: 53

